\documentclass[11pt,a4paper]{article}

\usepackage[utf8]{inputenc}
\usepackage[T1]{fontenc}
\usepackage{amsmath,amssymb,amsfonts}
\usepackage{physics}
\usepackage{graphicx}
\usepackage{booktabs}
\usepackage{array}
\usepackage{siunitx}
\usepackage{hyperref}
\usepackage[margin=1in]{geometry}
\usepackage{xcolor}
\usepackage{listings}
\usepackage{enumitem}
\usepackage{fancyhdr}
\usepackage{titlesec}

\hypersetup{
    colorlinks=true,
    linkcolor=blue!60!black,
    urlcolor=blue!60!black,
    citecolor=blue!60!black,
    pdftitle={PY-BEMCS Physics and User Manual},
    pdfauthor={Dr. Bharat Singh Rawat}
}

\pagestyle{fancy}
\fancyhf{}
\fancyhead[L]{PY-BEMCS Manual}
\fancyhead[R]{\thepage}
\fancyfoot[C]{\small Python Beam Extraction \& Monte Carlo Simulator}

\titleformat{\section}{\Large\bfseries\color{blue!50!black}}{\thesection}{1em}{}
\titleformat{\subsection}{\large\bfseries\color{blue!40!black}}{\thesubsection}{1em}{}

\newcommand{\PYBEMCS}{PY-BEMCS}

\begin{document}

% ---------------------- TITLE ----------------------
\begin{titlepage}
\centering
\vspace*{3cm}

{\Huge\bfseries Python Beam Extraction \&\\[0.3em] Monte Carlo Simulator \par}
\vspace{0.5cm}
{\Large\bfseries (\PYBEMCS)\par}

\vspace{1.5cm}
{\large\textit{Physics Reference and User Manual}\par}

\vspace{3cm}
{\large Dr.~Bharat Singh Rawat\par}
\vspace{0.3cm}
{\normalsize \href{mailto:bharat.bharat22@gmail.com}{bharat.bharat22@gmail.com}\par}

\vfill
{\normalsize \today\par}
\end{titlepage}

% ---------------------- ABSTRACT / MOTIVATION ----------------------
\begin{center}
{\Large\bfseries Abstract}
\end{center}
\vspace{0.5em}

The primary inspiration behind \PYBEMCS{} originated during my doctoral
research, where I worked extensively with ion sources and beam
extraction systems. Throughout that time, I evaluated numerous
simulation tools that utilised the Finite Element Method (FEM) to model
beam extraction processes. However, the vast majority of these were
proprietary commercial codes locked behind prohibitive license fees and
recurring service charges.

My vision for \PYBEMCS{} is to break down these barriers by providing a
powerful, fully open-source beam extraction simulator dedicated to
non-commercial research by academics and students.  A core design
philosophy of this project is \emph{accessibility}: it is engineered to
be as user-friendly as possible, removing the steep learning curve of
complex programming so that researchers can focus directly on the
physics.  As detailed in this document, the resulting engine uses a
hybrid Particle-in-Cell (PIC) / fluid approach to model plasma
dynamics, coupled with Monte Carlo Collisions (MCC) for neutral
interactions, and a finite-difference thermal model augmented with a
cantilever-beam deformation model for structural degradation.

I officially initiated this project on my 34th birthday
(26/03/2026), as a personal commitment to fostering accessible, open
science for the next generation of researchers.  This document details
the physical models, numerical methods, and multiphysics couplings
implemented in \PYBEMCS{}, including the recently added configurable
beam species, user-imported reaction cross-sections, multi-material
grid library, and thermal cantilever deformation.

\vspace{1em}
\hrule
\vspace{1em}

\tableofcontents
\newpage

% =====================================================
\section{Introduction}
% =====================================================

\PYBEMCS{} is a 2D3V (two spatial, three velocity components)
electrostatic Particle-in-Cell (PIC) simulator for ion-beam extraction,
charge-exchange (CEX) plume physics, grid erosion, and coupled
thermo-mechanical response in gridded ion thrusters and similar
ion sources.

The code is written in Python using the Taichi framework for parallel
GPU/CPU kernels, with a PyQt5 graphical front-end.  A companion C++
3D Particle-in-Cell framework (\texttt{Cpp3D}) is also included.
This manual documents the physics model used by the Python engine
(\texttt{physics\_engine.py}).

\subsection{Coordinate system and mesh}

A Cartesian (\(x,y\)) mesh is used, where \(x\) is the beam
axial direction and \(y\) is the radial direction (axisymmetric
interpretation).  The computational domain has dimensions
\(L_x \times L_y\) with cell size \(\Delta x \times \Delta y\):
\begin{equation}
    n_x = \frac{L_x}{\Delta x} + 1,\qquad
    n_y = \frac{L_y}{\Delta y} + 1.
\end{equation}
Particles carry three velocity components (\(v_x, v_y, v_z\)) with
\(v_z\) conserved by symmetry during the push but free to be modified
by collisions.

% =====================================================
\section{Particle-in-Cell Algorithm}
% =====================================================

A single time step of the explicit electrostatic PIC loop is:
\begin{enumerate}[nosep]
    \item Inject ions (and optionally electrons) at the source.
    \item Deposit charge on the mesh.
    \item Solve Poisson's equation for the potential.
    \item Compute the electric field.
    \item Push particles using the Boris algorithm.
    \item Apply collisions (CEX, SEE, ionisation).
    \item Handle wall hits, sputtering, thermal deposition.
    \item Remove dead particles.
\end{enumerate}
All grid operations are vectorised in NumPy and the computationally
expensive kernels (density deposition, particle push, thermal
diffusion) run through Taichi.

% =====================================================
\section{Electrostatic Field Model}
% =====================================================

\subsection{Poisson equation with Boltzmann electrons}

The self-consistent plasma potential satisfies
\begin{equation}
    \nabla^{2}\phi \;=\; -\,\frac{1}{\varepsilon_{0}}
        \left[ q_{\mathrm{i}}\,n_{\mathrm{i}}(\vb*{x})
             - e\,n_{\mathrm{e}}(\vb*{x}) \right],
    \label{eq:poisson}
\end{equation}
with \(n_{\mathrm i}\) obtained directly from the ion macro-particles
and \(n_{\mathrm e}\) from the inertia-less Boltzmann relation
\begin{equation}
    n_{\mathrm{e}}(\vb*{x}) \;=\; n_{0}\,
    \exp\!\left[\frac{\min(\phi,\phi_{\mathrm{pl}})-\phi_{\mathrm{pl}}}{T_{e}}\right],
    \label{eq:boltzmann}
\end{equation}
where \(T_{e}\) is in electron-volts, \(n_{0}\) is the upstream plasma
density, and \(\phi_{\mathrm{pl}}\) is the plasma potential applied at the
left boundary.  The \(\min()\) prevents unphysical over-density when the
local potential exceeds the plasma reference.

Equation~\eqref{eq:poisson} is discretised with a standard
five-point stencil
\begin{equation}
    \phi_{i+1,j}+\phi_{i-1,j}+\phi_{i,j+1}+\phi_{i,j-1}-4\phi_{i,j}
    = -\frac{(\Delta x)^{2}}{\varepsilon_{0}}\,\rho_{i,j},
\end{equation}
and inverted with a sparse LU factorisation performed once per
remesh.  The electron source term is updated in a successive
under-relaxation loop (relaxation factor \(\omega = 0.2\)).

\subsection{Charge deposition}

Macro-particles deposit their charge to the nearest mesh node with
bilinear weighting:
\begin{equation}
    \rho_{\mathrm{i},j} \;+\!=\; \frac{q_{\mathrm{i}}\,w}{V_{\mathrm{cell}}},
    \qquad
    V_{\mathrm{cell}} = \Delta x\,\Delta y\,\Delta z,
\end{equation}
where \(w\) is the macro-weight (physical particles per simulation
particle) and \(\Delta z = \SI{1}{\milli\meter}\) is the assumed
out-of-plane cell depth.

\subsection{Macro-particle weight}
\label{sec:injection-weight}

For a target of \(N_{\mathrm{ppc}}\) particles per cell we set
\begin{equation}
    w \;=\;
    \max\!\left(\frac{n_{0}\,V_{\mathrm{cell}}}{N_{\mathrm{ppc}}},\;10^{3}\right),
\end{equation}
with \(N_{\mathrm{ppc}} = 40\) by default.

\subsection{Electric field}

\begin{equation}
    \vb{E} \;=\; -\,\grad\phi,
    \qquad
    E_{x,i,j} = -\frac{\phi_{i+1,j}-\phi_{i-1,j}}{2\Delta x},\quad
    E_{y,i,j} = -\frac{\phi_{i,j+1}-\phi_{i,j-1}}{2\Delta y}.
\end{equation}

% =====================================================
\section{Particle Motion: 2D3V Boris Pusher}
% =====================================================

Particles are advanced in time with the Boris leap-frog scheme
\cite{birdsall}.  Starting from \(\vb{v}^{n-1/2}\) and positions
\(\vb{x}^{n}\):
\begin{align}
    \vb{v}^{-} &= \vb{v}^{n-1/2} + \frac{q\,\vb{E}}{2m}\,\Delta t,\\
    \vb{t} &= \frac{q\,\vb{B}}{2m}\,\Delta t, \qquad
    \vb{s} = \frac{2\,\vb{t}}{1 + \abs{\vb{t}}^{2}},\\
    \vb{v}' &= \vb{v}^{-} + \vb{v}^{-}\cross\vb{t},\\
    \vb{v}^{+} &= \vb{v}^{-} + \vb{v}'\cross\vb{s},\\
    \vb{v}^{n+1/2} &= \vb{v}^{+} + \frac{q\,\vb{E}}{2m}\,\Delta t,\\
    \vb{x}^{n+1} &= \vb{x}^{n} + \vb{v}^{n+1/2}\,\Delta t.
\end{align}
The algorithm is second-order accurate and conserves kinetic
energy exactly when \(\vb{E}=0\).  It is implemented as a Taichi
kernel operating in parallel over all particles.

The charge-to-mass ratio in the pusher uses the
user-configured ion mass and charge state:
\begin{equation}
    \frac{q}{m} = \frac{Z\,e}{m_{\mathrm{ion}}},
    \qquad m_{\mathrm{ion}} = A\,u,
\end{equation}
where \(A\) is the atomic mass in amu and \(Z\) is the charge state.

% =====================================================
\section{Ion Beam Extraction}
% =====================================================

\subsection{Bohm sheath criterion}

Ions are injected through the screen-grid hole with a drift
velocity equal to the Bohm velocity,
\begin{equation}
    v_{\mathrm{B}} \;=\; \sqrt{\frac{Z\,e\,T_{e}}{m_{\mathrm{ion}}}},
    \label{eq:bohm}
\end{equation}
yielding an injection current
\begin{equation}
    I_{\mathrm{i}} \;=\; 0.61\,Z\,e\,n_{0}\,v_{\mathrm{B}}\,A_{\mathrm{inj}},
\end{equation}
where \(A_{\mathrm{inj}}\) is the open area of the first grid and
the \(0.61\) factor is the classical pre-sheath correction.
The number of macro-particles injected per step is
\begin{equation}
    N_{\mathrm{inj}} \;=\; \frac{I_{\mathrm{i}}\,\Delta t}{Z\,e\,w},
\end{equation}
with a stochastic rounding step so the time-averaged current is
exactly \(I_{\mathrm{i}}\).

\subsection{Injection velocity distribution}

Each injected ion has
\begin{equation}
    v_{x} = v_{\mathrm{B}} + \xi_{x}\,v_{\mathrm{th}},\qquad
    v_{y,z} = \xi_{y,z}\,v_{\mathrm{th}},\qquad
    v_{\mathrm{th}} = \sqrt{\frac{Z\,e\,T_{\mathrm{i}}}{m_{\mathrm{ion}}}},
\end{equation}
where \(\xi_i\) are zero-mean unit-variance Gaussians and
\(T_{\mathrm i}\) is the ion temperature in eV.

% =====================================================
\section{Multi-Grid Optics and Geometry}
% =====================================================

The user specifies an arbitrary number of grids, each defined by
voltage \(V_{g}\), thickness \(t\), gap to next \(g\), hole radius
\(r\), and chamfer angle \(\theta\).  The grid mask is constructed
as
\begin{equation}
    \mathcal{M}_{g} \;=\;
    \left\{(x,y)\,\middle|\, x_{s}\leq x \leq x_{s}+t,\;
    y \geq r + (x-x_{s})\tan\theta\right\}.
\end{equation}

\subsection{Cantilever thermal deformation}

Each grid is treated as a cantilever beam clamped at the top wall
(\(y=L_{y}\)) and free at the aperture edge (\(y=r\)).  For a
uniform temperature rise \(\Delta T\) of a thin plate of length
\(L=L_{y}-r\) and thickness \(t\) the maximum tip deflection is
\begin{equation}
    \delta_{\max} \;=\; \frac{\alpha\,\Delta T\,L^{2}}{2\,t},
    \label{eq:cantilever}
\end{equation}
with \(\alpha\) the linear thermal-expansion coefficient.  The
deflection profile follows the Euler--Bernoulli cantilever shape
\begin{equation}
    \delta(y) \;=\; \delta_{\max}\,\eta^{2},
    \qquad \eta = \frac{L_{y}-y}{L}\in[0,1],
\end{equation}
and is applied as an \(x\)-offset to the grid boundary mask,
producing visible bowing of the grid toward or away from the beam
axis.  A remesh is triggered whenever the deflection changes by
more than $5\,\mu\mathrm{m}$.

\subsection{RF co-extraction}

The voltage on a user-selected grid can be modulated at RF
frequency~\(f\):
\begin{equation}
    V_{g}(t) \;=\; V_{\mathrm{DC}} + V_{\mathrm{RF}}\sin(2\pi f t).
\end{equation}
The Poisson equation is re-solved every step when RF is enabled.
With co-extraction enabled, the source also emits electrons at the
Bohm flux using the saturation current
\(I_{e}=0.25\,e\,n_{0}\,\bar{v}_{e}\,A_{\mathrm{inj}}\).

% =====================================================
\section{Collisional Processes}
% =====================================================

\subsection{Charge-exchange collisions}

The local neutral density follows Roy's analytical plume model
\begin{equation}
    n_{\mathrm{n}}(r,z) \;=\; n_{n,0}\,a\,
    \left[1 - \frac{1}{\sqrt{1+(r_{T}/R)^{2}}}\right]
    \cos\theta,
    \label{eq:roy}
\end{equation}
with \(a=(1-1/\sqrt{2})^{-1}\),
\(R=\sqrt{r^{2}+(z+r_{T})^{2}}\),
\(\theta=\arctan(r/(z+r_{T}))\), and \(r_{T}\) the maximum
transverse domain size.

The collision probability for each primary ion over time \(\Delta t\)
follows the null-collision Monte Carlo Method (MCC) of
Birdsall~\cite{birdsall}:
\begin{equation}
    P_{\mathrm{CEX}} \;=\; 1-\exp\!\left[-n_{\mathrm{n}}\,
       \sigma_{\mathrm{CEX}}(g)\,g\,\Delta t\right],
    \label{eq:mcc}
\end{equation}
where \(g=\abs{\vb{v}}\) is the relative speed.  The default
analytical cross section is the empirical Xe\(^{+}\)--Xe formula
\begin{equation}
    \sigma_{\mathrm{CEX}}(g) \;=\;
    \left[-0.8821\ln g + 15.1262\right]^{2}\times 10^{-20}\;\si{\meter\squared}.
    \label{eq:sigma_cex}
\end{equation}
If the user has imported a CX cross-section table,
equation~\eqref{eq:sigma_cex} is replaced by a cubic spline
interpolation of the tabulated \(\sigma(E)\) data in log--log
space (see Sec.~\ref{sec:csimport}).

When a collision occurs the ion velocity is replaced by a
Maxwellian sample at the neutral temperature using Birdsall's
\(M\!=\!3\) acceptance scheme:
\begin{equation}
    v_{i}' \;=\; v_{\mathrm{th},n}\,f_{M}, \qquad
    f_{M} = 2\left(R_{1}+R_{2}+R_{3}-\tfrac{3}{2}\right),
\end{equation}
with \(R_{i}\sim U(0,1)\) and
\(v_{\mathrm{th},n}=\sqrt{2k_{B}T_{n}/m_{\mathrm{ion}}}\).

\subsection{Secondary electron emission (SEE)}

Ions striking a grid surface with energy \(E_{\mathrm{eV}}\)
release secondary electrons with probability
\begin{equation}
    \gamma(E) \;=\; \gamma_{0}+\beta\,E_{\mathrm{eV}},
    \qquad \gamma_{0}=0.05,\;\beta=10^{-4},
\end{equation}
clipped to \([0,1]\).  When a user-supplied SEE yield table is
loaded, \(\gamma(E)\) is replaced by a spline of the imported
data.  Emitted electrons have a half-Maxwellian velocity
distribution at \(T_{\mathrm{SEE}}=\SI{2}{\electronvolt}\).

\subsection{User-imported cross-section data}
\label{sec:csimport}

The GUI \textbf{Beam $\rightarrow$ Cross-Section Manager} accepts
CSV tables of energy \([\si{\electronvolt}]\) versus cross-section
\([\si{\meter\squared}]\) for CX, SEE, or custom reactions.
A cubic spline \(s(\log_{10}E)\) is fitted in log--log space with
user-selectable smoothing factor:
\begin{equation}
    \log_{10}\sigma(E) \;\approx\; s\!\left(\log_{10}E\right).
\end{equation}
During simulation the spline is evaluated at each particle
energy, with clamping to the tabulated data range.  Multiple
datasets can be loaded simultaneously and plotted.

% =====================================================
\section{Sputtering and Erosion}
% =====================================================

\subsection{Yield and damage accumulation}

Each ion impact on a grid boundary contributes to an accumulated
damage map
\begin{equation}
    D_{i,j} \;+\!=\; Y(E)\,\kappa_{\mathrm{acc}},
    \label{eq:damage-accum}
\end{equation}
with the energy-dependent yield (empirical Matsunami fit)
\begin{equation}
    Y(E) \;=\;
    \min\!\left[
    Y_{0}\,\bigl(E_{\mathrm{eV}}-E_{\mathrm{th}}\bigr)^{1.5},\;1
    \right],
\end{equation}
where \(Y_{0}\) is the material-dependent yield coefficient
(atoms per incident ion per \si{\eV\tothe{1.5}}) and
\(E_{\mathrm{th}}\) is the sputter threshold energy.
\(Y(E)\) is the dimensionless number of target atoms
ejected per incident ion.

\subsection{Time-acceleration factor \texorpdfstring{$\kappa_{\mathrm{acc}}$}{kappa\_acc}}
\label{sec:kappa-acc}

Real grid erosion in an ion thruster occurs over
\(10^{3}\)--\(10^{4}\,\si{\hour}\) of firing, whereas a PIC
simulation can only afford \(\mathcal{O}(\si{\milli\second})\)
of physical time. \(\kappa_{\mathrm{acc}}\) is a
dimensionless time-compression factor that multiplies the
per-hit damage in Eq.~\eqref{eq:damage-accum} so that a short
simulation run reproduces the total eroded material of a long
real-world firing.

Formally, let \(T_{\mathrm{sim}}\) be the total physical
time simulated and \(T_{\mathrm{real}}\) be the operating
duration being emulated. Because the incident ion flux is
stationary once the beam reaches steady state, the total
number of impacts scales linearly with time, and so does
the accumulated damage. The consistent choice is therefore
\begin{equation}
    \kappa_{\mathrm{acc}} \;=\; \frac{T_{\mathrm{real}}}{T_{\mathrm{sim}}},
    \label{eq:kappa-acc}
\end{equation}
so that \(1\,\si{\milli\second}\) of simulated firing
represents \(\kappa_{\mathrm{acc}}\,\si{\milli\second}\) of
real thruster operation. For example, compressing
\(T_{\mathrm{real}}=1000\,\si{\hour}\) of operation into
\(T_{\mathrm{sim}}=10\,\si{\milli\second}\) of simulation
requires \(\kappa_{\mathrm{acc}}=3.6\times10^{8}\).

\paragraph{Calibration against measured erosion rate.}
\(\kappa_{\mathrm{acc}}\) can also be obtained by matching
the simulated beginning-of-life erosion rate against a
published experimental value. If \(\dot m_{\mathrm{meas}}\)
(\si{\kilo\gram\per\second}) is the measured mass-loss rate
of a reference thruster and the simulation removes a mass
\(m_{\mathrm{sim}}\) during \(T_{\mathrm{sim}}\) with
\(\kappa_{\mathrm{acc}}=1\), then
\begin{equation}
    \kappa_{\mathrm{acc}} \;=\;
    \frac{\dot m_{\mathrm{meas}}\,T_{\mathrm{sim}}}
         {m_{\mathrm{sim}}},
\end{equation}
which forces the simulated erosion rate to match
measurement at \(t=0\). The two definitions are equivalent
when the yield model and ion flux are correctly calibrated.

\subsection{Damage threshold \texorpdfstring{$D_{\mathrm{fail}}$}{D\_fail}}
\label{sec:d-fail}

When the accumulated damage in a cell exceeds
\(D_{\mathrm{fail}}\), the cell is marked non-solid, the
binary boundary mask \(\mathrm{isBound}\) is cleared, and
the sparse Poisson operator is refactorised so that the
following field solve sees the updated geometry. Physically
this represents the moment at which enough material has
been sputtered away from that cell that it no longer
contributes to the electrostatic boundary.

The dimensionless \(D_{\mathrm{fail}}\) can be mapped to a
physical wall-volume threshold as follows. A single ion
removes \(Y(E)\) atoms, each occupying an atomic volume
\begin{equation}
    V_{\mathrm{atom}} \;=\;
    \frac{M_{\mathrm{target}}}{\rho_{\mathrm{target}}\,N_{A}},
\end{equation}
where \(M_{\mathrm{target}}\) is the molar mass and
\(\rho_{\mathrm{target}}\) the density of the grid
material (e.g.\ \(V_{\mathrm{atom}}\approx 1.55\times10^{-29}\,
\si{\meter\cubed}\) for molybdenum). A macroparticle impact
therefore sputters a real volume
\begin{equation}
    \Delta V_{\mathrm{hit}} \;=\;
    Y(E)\,w\,V_{\mathrm{atom}}\,\kappa_{\mathrm{acc}},
\end{equation}
with \(w\) the macroparticle weight defined in
Sec.~\ref{sec:injection-weight}. A cell fails when its
accumulated \(\Delta V\) equals the physical material
volume present at that face,
\begin{equation}
    V_{\mathrm{cell}} \;=\;
    \Delta x\,\Delta y\,t_{\mathrm{wall}},
\end{equation}
where \(t_{\mathrm{wall}}\) is the local wall thickness
(for the accelerator grid, typically \(t_{a}\) from the
geometry inputs). The physically motivated threshold is
therefore
\begin{equation}
    D_{\mathrm{fail}} \;=\;
    \frac{V_{\mathrm{cell}}}{w\,V_{\mathrm{atom}}}
    \;=\;
    \frac{\Delta x\,\Delta y\,t_{\mathrm{wall}}}
         {w\,V_{\mathrm{atom}}},
    \label{eq:d-fail-physical}
\end{equation}
which is the number of sputter events (weighted by yield)
required to hollow out one cell of grid material.
Eq.~\eqref{eq:d-fail-physical} converts an otherwise
arbitrary GUI knob into a quantity that depends only on
mesh resolution, macroparticle weight, and target material
properties.

% =====================================================
\section{Thermal Model}
% =====================================================

\subsection{Energy deposition}

The kinetic energy of every ion absorbed by a grid cell is
converted to heat with an acceleration factor
\(\kappa_{\mathrm{th}}\) that compresses the
diffusive time-scale:
\begin{equation}
    \Delta T_{\mathrm{heat}} \;=\;
    \frac{\tfrac{1}{2}m_{\mathrm{ion}} v^{2}\,w}{C_{\mathrm{cell}}}\,
    \kappa_{\mathrm{th}},
\end{equation}
with cell heat capacity
\(C_{\mathrm{cell}}=\rho\,\Delta x\,\Delta y\,\Delta z\,c_{p}\).

\subsection{Radiative cooling}

Each boundary cell loses energy by Stefan--Boltzmann radiation:
\begin{equation}
    \Delta T_{\mathrm{rad}} \;=\;
    \frac{\varepsilon\,\sigma_{\mathrm{SB}}\,A_{\mathrm{cell}}\,\Delta t\,
         \kappa_{\mathrm{th}}}{C_{\mathrm{cell}}}
    \left(T^{4}-T_{\infty}^{4}\right),
\end{equation}
with \(T_{\infty}=\SI{300}{\kelvin}\),
\(\sigma_{\mathrm{SB}}=5.67\times10^{-8}\,\si{\watt\per\meter\squared\per\kelvin\tothe{4}}\),
and \(\varepsilon\) the material emissivity.

\subsection{2D thermal conduction}

An explicit five-point FTCS stencil is used on the solid mask
only,
\begin{equation}
    T_{i,j}^{n+1} = T_{i,j}^{n}
    + \mathrm{Fo}_{x}\,(T_{i-1,j}-2T_{i,j}+T_{i+1,j})
    + \mathrm{Fo}_{y}\,(T_{i,j-1}-2T_{i,j}+T_{i,j+1}),
\end{equation}
with diffusivity
\(\alpha_{\mathrm{diff}}=k/(\rho c_{p})\) and Fourier numbers
\(\mathrm{Fo}=\alpha_{\mathrm{diff}}\,\Delta t\,\kappa_{\mathrm{th}}/(\Delta x)^{2}\)
automatically rescaled to satisfy the stability bound
\(\mathrm{Fo}\leq 0.2\).  Insulating Neumann conditions are applied
wherever a neighbouring cell is not solid.

\subsection{Material properties}

Eight grid materials are supported through the
\textbf{Materials $\rightarrow$ Grid Material} dialog.
All properties are editable; the presets are:

\begin{center}
\small
\begin{tabular}{lccccccc}
\toprule
Material & \(k\) & \(\rho\) & \(c_{p}\) & \(\varepsilon\) &
    \(\alpha\) & \(E\) & \(Y_{0}\)\\
         & (W/mK) & (kg/m\(^3\)) & (J/kgK) & & (1/K) & (GPa) & \\
\midrule
Molybdenum    & 138 & 10\,280 & 250 & 0.80 & 4.8e-6  & 329 & 1.05e-4\\
Steel SS316   & 16.3& 8\,000  & 500 & 0.60 & 16.0e-6 & 193 & 2.8e-4\\
Titanium      & 21.9& 4\,507  & 520 & 0.50 & 8.6e-6  & 116 & 1.8e-4\\
Graphite      & 120 & 2\,200  & 710 & 0.85 & 3.0e-6  & 11  & 3.5e-4\\
\bottomrule
\end{tabular}
\end{center}

% =====================================================
\section{Neutralizer and Co-Extraction}
% =====================================================

A localised electron source at
\((x_{\mathrm{n}},y_{\mathrm{n}})\) injects a user-set number
\(\dot N_{e}\) of macro-electrons per step with a thermal
Maxwell--Boltzmann distribution at temperature \(T_{e}\):
\begin{equation}
    v_{x,y,z} \;=\; \xi \sqrt{\tfrac{2 e T_{e}}{m_{e}}},
    \qquad \xi\sim\mathcal N(0,1).
\end{equation}
The reduced electron mass
\(m_{e}=m_{\mathrm{ion}}/\mathcal{R}\) with
\(\mathcal{R}\sim100\)--\(10^{3}\) is an \emph{artificial}
mass ratio used for speed-up; it scales the gyro-frequency
and plasma frequency by \(1/\sqrt{\mathcal R}\) but preserves
the qualitative plume structure.

% =====================================================
\section{Beam Divergence and Diagnostics}
% =====================================================

The beam divergence is computed from the 95\textsuperscript{th}
percentile of the angular distribution of primary ions past the
last grid:
\begin{equation}
    \theta_{95} \;=\;
    Q_{0.95}\!\left[\,\abs{\arctan\!\left(v_{y}/v_{x}\right)}\,\right].
\end{equation}
The minimum centerline potential (electron backstreaming
saddle) is sampled at the middle of the second grid:
\begin{equation}
    V_{\mathrm{sp}} \;=\; \phi\!\left(y=0,\,x=x_{g2,\mathrm{mid}}\right).
\end{equation}
Both quantities are logged every five iterations and can be
exported to CSV, together with the instantaneous grid
temperatures.

% =====================================================
\section{Numerical parameters and stability}
% =====================================================

\begin{itemize}[nosep]
    \item Time step: \(\Delta t = \SI{1}{\nano\second}\).  For
        Xe\(^{+}\) at \SI{1}{\kilo\electronvolt},
        \(v_{i}\approx 4\times10^{4}\,\si{\meter\per\second}\),
        giving \(v_{i}\Delta t\ll\Delta x\).
    \item Mesh spacing: \(\Delta x = \Delta y = 15~\mu\mathrm{m}\).
    \item Debye length: \(\lambda_{D}=\sqrt{\varepsilon_{0}T_{e}/(e n_{0})}\)
        is \(\gtrsim\Delta x\) for
        \(n_{0}\leq 10^{17}\,\si{\per\meter\cubed}\) and
        \(T_{e}\geq\SI{1}{\electronvolt}\).
    \item Poisson iterations per push: 2 ordinary,
        30 at domain build.
    \item Thermal acceleration: \(\kappa_{\mathrm{th}}=10^{7}\)
        (maps microsecond PIC timescale to realistic seconds-long
        heating).
\end{itemize}

% =====================================================
\section{GUI Workflow Summary}
% =====================================================

\begin{enumerate}[nosep]
    \item \textbf{Beam $\rightarrow$ Ion Species:}
        choose propellant (Xe, Kr, Ar, N\(_2\), O\(_2\), H\(_2\),
        He, Hg, Cs) or set custom mass and charge state.
    \item \textbf{Beam $\rightarrow$ Cross-Section Manager:}
        import CSV data for CX, SEE, or custom reactions and fit
        splines.
    \item \textbf{Materials $\rightarrow$ Grid Material:}
        pick material preset or define custom thermophysical
        properties.
    \item Configure grid geometry, plasma, RF, and neutralizer
        in the main control panel.
    \item Press \textbf{Build Domain}, then \textbf{Start Beam}.
    \item Monitor live diagnostics; export telemetry or particle
        kinematics to CSV; optionally record animation GIF.
\end{enumerate}

% =====================================================
\section{Headless Operation}
% =====================================================

The same physics engine can be driven from a configuration file:
\begin{verbatim}
python Python/run_simulation_from_config.py
\end{verbatim}
All parameters (beam species, grid material, cross-section
files, grids, plasma, RF, neutralizer, terminal output) are read
from \texttt{Python/config.ini}.  This mode is intended for
scripted parameter sweeps and batch runs.

% =====================================================
\begin{thebibliography}{9}
\bibitem{birdsall}
C.~K. Birdsall and A.~B. Langdon,
\textit{Plasma Physics via Computer Simulation},
Institute of Physics Publishing, 1991.

\bibitem{roy}
R.~I.~S. Roy, \textit{Numerical Simulation of Ion Thruster Plume
Backflow}, Ph.D. Thesis, MIT, 1995.

\bibitem{matsunami}
N.~Matsunami \textit{et al.},
``Energy dependence of the yields of ion-induced sputtering of
monatomic solids,'' \textit{Atomic Data and Nuclear Data Tables},
31, 1--80 (1984).
\end{thebibliography}

\end{document}
